\chapter{Comprendre la limite d’une suite}
\section*{Exercice 9.1 -- Définition et rang}
$3/(n+1)<1/100$ équivaut à $n+1>300$, donc $n\ge300$. Le rang $299$ donne l'égalité $0{,}01$ et ne convient pas pour l'inégalité stricte.

\section*{Exercice 9.2 -- Bornée mais divergente}
La sous-suite paire tend vers $1$ et l'impaire vers $-1$. Si la suite entière avait une limite, toutes ses sous-suites auraient cette même limite, par la définition « à partir d'un rang ». Les deux limites différentes excluent la convergence.

\section*{Exercice 9.3 -- Gendarmes}
$-1/n\le\sin n/n\le1/n$. Les deux bornes tendent vers zéro, donc la suite tend vers zéro. Il n'est pas nécessaire que $\sin n$ converge.

\section*{Exercice 9.4 -- Non majorée}
Poser $u_{2n}=2n+1$ et $u_{2n+1}=1$. La suite n'est pas majorée, mais revient à $1$ à des indices arbitrairement grands ; le seuil $2$ n'est donc jamais dépassé durablement.

\section*{Exercice 9.5 -- Ordre et limite}
Prendre $\varepsilon=\ell/2$. Pour les grands indices, $|u_n-\ell|<\ell/2$, donc $u_n>\ell/2>0$. Cela n'impose aucun signe aux premiers termes.

\section*{Exercice 9.6 -- Négation de convergence}
Il existe $\varepsilon>0$ tel que, pour tout $N$, il existe $n\ge N$ avec $|u_n|\ge\varepsilon$. Le même $\varepsilon$ doit convenir à tous les rangs $N$. Dire seulement que des termes sont non nuls ne suffit pas.
